\section{Institutional Background}
\label{sec:institutional}

\subsection{Brazilian Procurement Law and the Convite Modality}
\label{sec:lei8666}

Brazil's general procurement statute, Lei 8.666/1993, establishes
five modalities for public purchasing, ranked by contract value:
\emph{convite} (lowest threshold), \emph{tomada de pre\c{c}os}
(intermediate), \emph{concorr\^encia} (highest), plus two
special-purpose formats (\emph{concurso} for technical contests and
\emph{leil\~ao} for asset sales). The modality thresholds during
most of our sample period (2009--2018) were R\$80,000 for goods
and services under convite and R\$150,000 for engineering works;
Decreto 9.412/2018 raised these to R\$176,000 and R\$330,000,
respectively, affecting only the last year of our data.%
\footnote{The original 1993 thresholds were denominated in
\emph{cruzeiros reais} and updated periodically by executive
decree. All figures in this paper use nominal values.}

Convite is the simplest and least formal modality. The procuring
unit selects and invites at least three firms by letter
(\emph{carta-convite}); other registered vendors may
self-nominate up to 24 hours before the submission deadline
(Art.\ 22, \S\,3). Proposals are submitted in sealed envelopes and
opened in a public session; the lowest qualifying price wins.
There is no live bidding: convite is a first-price, sealed-bid
auction.

\subsection{The Minimum-Bidder Rule}
\label{sec:min_bidder}

The institutional feature at the center of this paper is
Art.\ 22, \S\,3 of Lei 8.666/1993, which requires the procuring
unit to invite ``no fewer than 3 (three)'' participants in
convite tenders. Complementing this, Art.\ 22, \S\,7 specifies
the consequence when the floor is not met:

\begin{quote}
\small
When, due to market limitations or manifest lack of interest from
the invited parties, it is impossible to obtain the minimum number
of bidders required in \S\,3 of this article, these circumstances
must be duly justified in the process file, under penalty of
repetition of the convite.%
\footnote{Our translation. Original: ``Quando, por limita\c{c}\~oes
do mercado ou manifesto desinteresse dos convidados, for
imposs\'ivel a obten\c{c}\~ao do n\'umero m\'inimo de licitantes
exigidos no \S\,3\textordmasculine{} deste artigo, essas
circunst\^ancias dever\~ao ser devidamente justificadas no
processo, sob pena de repeti\c{c}\~ao do convite.'' (Art.\ 22,
\S\,7, Lei 8.666/1993.)}
\end{quote}

The practical implication is sharp. A procuring unit that receives
fewer than three bids faces two options: (i)~document that the
market cannot supply three competitors---a written justification
subject to audit by the Tribunal de Contas---or (ii)~cancel and
re-advertise, which delays procurement and may attract unwanted
competition. Both options are costly for the procuring unit and
undesirable for a cartel. A cartel controlling the winning bid
therefore has a direct incentive to ensure that the three-bidder
threshold is met, which it does by deploying affiliated firms to
submit deliberately losing bids---\emph{cover bidders}.

The rule's logic is pro-competitive: more participants should
generate more price competition. But this logic assumes that
additional participants are genuine competitors. When a cartel
can manufacture participation at low cost---electronic bidding on
the BEC platform requires only CAUFESP registration and minimal
documentation---the rule's binding constraint falls not on
competition but on the cartel's organizational capacity. The
minimum-bidder floor becomes a mandate to produce synthetic
competition.

\subsection{Preg\~ao: The Alternative Regime}
\label{sec:pregao}

Lei 10.520/2002 introduced \emph{preg\~ao}, a reverse-auction
modality for ``common goods and services''---those whose quality
standards can be objectively specified by market reference.
Preg\~ao was designed to address the rigidity and formalism of
Lei 8.666 procedures: it has no value ceiling, allows live
price competition, and was mandated as the preferred modality for
routine procurement across the federal government (Decreto
5.450/2005).%
\footnote{For S\~ao Paulo state, Decreto 49.722/2005 extended the
BEC platform to include preg\~ao eletr\^onico.}

The mechanics differ from convite in two respects relevant to this
paper. \emph{First}, preg\~ao is an electronic reverse auction:
the lowest initial proposals are admitted to a live descending-bid
phase, and the session ends when no further reductions are
offered. The real-time format makes coordination harder but not
impossible; cover bidders can still inflate the initial price
envelope. \emph{Second}, and critically, preg\~ao has no
minimum-bidder requirement. A tender may proceed with a single
bidder. There is no statutory floor to satisfy, no justification
burden, and no cancellation risk if participation is thin.

This institutional difference generates the paper's central
empirical comparison. If the minimum-bidder rule facilitates
cover bidding, then cover-bidder presence should have a
qualitatively different relationship with prices across the two
regimes. In convite, the rule forces deployment even in markets
where cover bidding would otherwise be unprofitable (a corner
solution); in preg\~ao, deployment is purely voluntary (an
interior solution). We exploit this variation in
Section~\ref{sec:empirical}.

\subsection{The BEC Platform}
\label{sec:bec_platform}

The Bolsa Eletr\^onica de Compras (BEC/SP) is S\~ao Paulo
state's centralized electronic procurement platform, established
by Decreto 45.085/2000 and formally designated as BEC/SP by
Decreto 45.695/2001. Administered by the Coordenadoria de
Compras Eletr\^onicas of the Secretaria da Fazenda e Planejamento,
the platform serves all organs of the state's direct and indirect
administration, plus municipalities that have signed cooperation
agreements (\emph{conv\^enios}). During our sample period
(2009--2019), the platform encompasses 1,308 purchasing units
(PBUs) and handles both convite (sealed-bid) and preg\~ao
eletr\^onico (reverse auction) in fully electronic format.

Three features of BEC make it particularly suited to this study.
\emph{First}, both modalities operate on the same platform, so
the convite--preg\~ao comparison is not confounded by differences
in electronic infrastructure or vendor registration requirements.
All vendors register through CAUFESP (the state's unified supplier
registry); registration is free. \emph{Second}, BEC records all
bids, including losing bids, at the item level---a granularity
rare in procurement datasets and essential for identifying
participation patterns. \emph{Third}, the platform's electronic
nature means that bidding costs are low (staff time for bid
preparation and documentation, estimated at 1--2 hours per tender),
which sharpens the puzzle: if bidding is cheap and winning is
impossible, why do frequent losers persist?

\subsection{Cross-Country Context}
\label{sec:crosscountry}

Minimum-bidder provisions are not unique to Brazil.
Table~\ref{tab:crosscountry} summarizes participation floors across
six jurisdictions. The European Union's Directive 2014/24/EU
requires a minimum of five invited candidates in restricted
procedures and three in competitive dialogue, though the floors
apply to the invited set, not to bids received, and contracting
authorities may proceed with fewer if the market cannot supply more
(Art.~65). India's General Financial Rules (GFR 2017, Rule 162)
require that limited tender enquiries contact at least three
suppliers---a close structural parallel to Brazil's convite, with
similar value caps (Rs.~50 lakhs, approximately R\$300,000).
The United States' Federal Acquisition Regulation (FAR 14.101)
requires a ``reasonable expectation of receiving more than one
sealed bid'' before selecting the sealed-bid method, but allows
single-bid awards if competition was genuinely solicited. Neither
the WTO Government Procurement Agreement nor South Africa's PPPFA
imposes a numerical participation floor.

The common thread is that participation floors are most prevalent
in lower-value, invitation-based modalities---precisely the settings
where genuine competition is thinnest and the marginal cost of
manufacturing a cover bid is lowest. Brazil's convite is a
canonical case, but the mechanism we document---a participation
mandate that cartels satisfy by deploying affiliated
firms---generalizes to any jurisdiction with similar rules.

\begin{table}[htbp]
\centering
\caption{Minimum-Bidder Rules Across Jurisdictions}
\label{tab:crosscountry}
\begin{adjustbox}{max width=\textwidth}
\begin{threeparttable}
\small
\begin{tabular}{llcll}
\toprule
Jurisdiction & Rule & Floor & Scope & Consequence if unmet \\
\midrule
Brazil (Lei 8.666) & Art.\ 22, \S\S\,3+7 & 3 bids & Convite (low value) & Re-run or justify in writing \\
EU Dir.\ 2014/24 & Art.\ 65 & 3--5 invited & Restricted/competitive & May proceed if no others qualify \\
India (GFR 2017) & Rule 162 & 3+ contacted & Limited tender ($\leq$ Rs.\ 50L) & Implied re-solicitation \\
USA (FAR) & 14.101 & 2+ expected & Sealed bidding (method choice) & Method selection only \\
WTO GPA & Art.\ IX & None & All covered procurement & --- \\
South Africa (PPPFA) & 2022 Regs & None & All thresholds & --- \\
Brazil (Lei 14.133) & Art.\ 28 & None & All modalities & Convite abolished \\
\bottomrule
\end{tabular}
\begin{tablenotes}
\small
\item \textit{Notes:} ``Floor'' refers to the statutory minimum
on bidders (or invited candidates) required for the tender to
proceed. The EU floors apply to invited candidates, not received
bids. India's rule applies to contacted suppliers. The U.S.\ FAR
standard is prospective (about choosing the sealed-bid method) and
does not prevent single-bid awards. Brazil is not a party to the
WTO GPA; the comparison is illustrative.
\end{tablenotes}
\end{threeparttable}
\end{adjustbox}
\end{table}

\subsection{The 2021 Reform}
\label{sec:lei14133}

Lei 14.133/2021, Brazil's ``New Procurement Law,'' replaced
Lei 8.666 after a two-year transition period ending December 30,
2023. The new law abolishes the convite modality entirely
(Art.~28 retains only preg\~ao, concorr\^encia, concurso,
leil\~ao, and the new \emph{di\'alogo competitivo}). Modality
choice is now determined by the nature of the contract, not its
value. Because convite no longer exists, the three-bidder rule
no longer applies.%
\footnote{Lei 14.133 retains minimum qualification thresholds
for competitive dialogue and restricted procedures, but these
apply to vendor qualifications, not to the number of bids
received.}

This reform is not directly relevant to our empirical window
(2009--2019), during which both Lei 8.666 and Lei 10.520 were
fully operative. It does, however, provide a natural endpoint for
the institutional arrangement we study: the rule that created
demand for cover bidders has been legislated away. Whether the
abolition of convite has reduced cover-bidder activity in
practice is an open question that future research, using
post-2024 BEC data, could address.
